\apendice{Anexo de sostenibilización curricular}

\section{Introducción}
El desarrollo de este Trabajo de Fin de Grado (\textit{TFG}) ha supuesto una oportunidad inestimable para trascender la aplicación puramente técnica de los conocimientos adquiridos durante mi formación académica, permitiéndome integrar de forma transversal los principios fundamentales de la sostenibilidad. Como ingeniero, es imprescindible comprender que el diseño de \textit{software} y la arquitectura de sistemas de información no son disciplinas neutras; tienen un impacto directo y profundo sobre la sociedad, la economía y el medio ambiente. 

En consonancia con las directrices de la \textbf{\href{https://www.crue.org/wp-content/uploads/2020/02/Directrices_Sosteniblidad_Crue2012.pdf}{CRUE} (Conferencia de Rectores de las Universidades Españolas)} para la introducción de la sostenibilidad en el currículum universitario, este proyecto se ha concebido desde una visión holística. Se ha asumido la responsabilidad ética que conlleva el desarrollo de tecnologías aplicadas a un entorno tan crítico como la salud humana. A continuación, se detalla cómo esta plataforma analítica orientada a la \textit{URCCPQ} aborda y satisface las tres dimensiones clásicas del desarrollo sostenible: social, económica y ambiental.

\subsection{Dimensión Social: Ética, Privacidad y Calidad Asistencial}
La dimensión social constituye el pilar central sobre el que se sustenta la motivación de este proyecto. El propósito último del sistema es la transformación de datos brutos en conocimiento clínico accionable para la mejora continua de la atención al paciente crítico. Al automatizar y facilitar el cálculo de los indicadores de calidad y seguridad de la \textit{SEMICYUC}, la herramienta empodera al personal médico para auditar su propia práctica, identificar áreas de mejora y prevenir complicaciones graves (como neumonías asociadas a ventilación mecánica, úlceras por presión o infecciones por catéter). En última instancia, esta tecnología contribuye directamente a la protección del derecho fundamental a la vida, promoviendo una asistencia sanitaria más segura y equitativa.

Asimismo, el diseño del sistema respeta escrupulosamente los derechos de privacidad y la normativa vigente de protección de datos (\textit{LOPD}). La implementación de un control de acceso basado en roles (\textit{RBAC}), el procesamiento aislado en la \textit{intranet} hospitalaria (sin depender de nubes corporativas externas) y los algoritmos de anonimización para la extracción de datos destinados a la investigación, demuestran un fuerte compromiso ético. Se garantiza que la identidad de los pacientes permanezca protegida, mitigando los riesgos inherentes a la digitalización de historiales médicos. 

Por otro lado, el desarrollo de una interfaz \textit{web} intuitiva, operable desde cualquier navegador sin complejas instalaciones previas, promueve la inclusión y la accesibilidad universal para todos los facultativos, independientemente de sus competencias o habilidades tecnológicas previas.

\subsection{Dimensión Económica: Software Libre y Soberanía Tecnológica}
Desde una perspectiva de sostenibilidad económica, el proyecto ha apostado firmemente por la utilización de tecnologías de código abierto (\textit{Open Source}). El empleo de lenguajes y herramientas como \textit{Python}, \textit{Reflex}, \textit{PostgreSQL} y \textit{Docker} elimina las barreras de entrada económicas asociadas a las costosas licencias de \textit{software} privativo (como ocurre con programas estadísticos de pago tipo \textit{SPSS}\footnote{\href{https://www.ibm.com/es-es/products/spss-statistics}{IBM SPSS Statistics}}) o a las suscripciones recurrentes en la nube. Esta decisión democratiza el acceso a herramientas analíticas avanzadas, permitiendo que hospitales públicos o unidades clínicas con presupuestos limitados puedan beneficiarse de la innovación sin comprometer sus recursos financieros.

Además, la arquitectura orquestada mediante contenedores garantiza una alta mantenibilidad y viabilidad a largo plazo. El sistema no depende de un proveedor tecnológico específico (\textit{vendor lock-in}\footnote{Situación en la que una empresa depende tanto de un proveedor tecnológico específico que cambiar a otro resulta extremadamente difícil o costoso.}), lo que asegura que el hospital pueda gestionar, replicar, modificar y actualizar la plataforma de forma autónoma. Esto fomenta la soberanía tecnológica de la institución, reduciendo la dependencia externa y los costes de mantenimiento.

\subsection{Dimensión Ambiental: Eficiencia Energética y Reducción de la Huella Ecológica}
Aunque con frecuencia se subestima el impacto ecológico del código informático, este proyecto ha integrado los principios de la computación verde (\textit{Green Computing}\footnote{Uso eficiente, sostenible y ecológico de los recursos informáticos}). La decisión de contenerizar la aplicación optimiza al máximo el uso de los recursos del servidor, evitando la necesidad de desplegar múltiples máquinas virtuales pesadas que consumen un exceso de memoria \textit{RAM} y capacidad de procesamiento.

Específicamente en su variante de uso autónomo y portátil, el despliegue del sistema sobre un dispositivo \textit{Raspberry Pi} representa un avance significativo en eficiencia energética. Mientras que un servidor tradicional o un ordenador de sobremesa en un hospital puede requerir un consumo de entre 200 y 400 vatios por hora, la \textit{Raspberry Pi} opera de manera completamente funcional con un consumo inferior a los 5 vatios. Esta drástica reducción de la demanda eléctrica minimiza la huella de carbono asociada al procesamiento de datos.

Por último, la digitalización integral de los informes visuales y de las métricas clínicas de la unidad impulsa la transición institucional hacia una metodología sin papel. La capacidad del sistema para generar automáticamente archivos \textit{PDF} estructurados y compartibles reduce de forma directa el consumo de recursos materiales (papel, tinta, almacenamiento físico) y la generación de residuos ofimáticos en la planta hospitalaria.

\subsection{Conclusión}
La conceptualización y el desarrollo de esta plataforma clínica han evidenciado que la ingeniería de \textit{software}, cuando se alinea con propósitos sociales, posee un potencial transformador excepcional. Este proyecto no se limita a resolver un problema de procesamiento de datos médicos; refleja un compromiso activo con una evolución tecnológica innovadora que es éticamente responsable, económicamente viable y respetuosa con el medio ambiente, respondiendo con rigor a las exigencias de la actual crisis climática y a las necesidades de la sociedad contemporánea.